% §11.4 Sample Complexity
%
% Derives how many independent samples (or rounds, or repetitions) are
% needed to drive a Bernoulli-shaped failure probability below a target
% $\delta$. The workhorse is the multiplicative Chernoff bound, stated
% as a theorem with an MGF + Markov proof sketch. Two applications:
% Miller-Rabin round count (a clean closed form because per-round
% $\fprate \le 1/4$) and BPP majority-vote amplification (the Chernoff
% calculation supplied at \Cref{thm:bpp-amplification}'s deferral
% point). Closes with the comparison to the variance-based interval of
% chapter 7 §7.2: Chebyshev-style polynomial decay versus Chernoff-style
% exponential decay on the same binomial estimator.
%
% Cross-reference resolution status:
%
% DEFINED here:
%   \label{sec:sample-complexity}       -- this section
%   \label{thm:chernoff}                -- multiplicative Chernoff bound
%
% RESOLVED (earlier chapters, drafted in Plans 1-7):
%   \Cref{sec:binomial-estimator}       -- chapter 7 §7.2
%   \Cref{prop:binomial-unbiased}       -- chapter 7 §7.2
%   \Cref{prop:binomial-variance}       -- chapter 7 §7.2
%   \Cref{thm:composition}              -- chapter 3 §3.3
%   \Cref{ax:element-indep}             -- chapter 4 §4.2 Axiom 1
%
% RESOLVED (chapter 11, earlier in this plan):
%   \Cref{sec:miller-rabin-formal}      -- §11.1
%   \Cref{thm:miller-rabin-formal}      -- §11.1
%   \Cref{def:bpp}                      -- §11.3
%   \Cref{thm:bpp-amplification}        -- §11.3 (proof deferred to here)

\section{Sample Complexity}
\label{sec:sample-complexity}

Every randomized algorithm of the preceding three sections has a knob
labelled $k$: the number of Miller-Rabin rounds, the number of BPP
repetitions for majority vote, the number of Karger trials. Each
$k$-knob does the same job: drive a per-trial error rate below a
target failure probability $\delta$. The sample-complexity question
is how to set the knob. Given a target relative error $\epsilon$ and
a target failure probability $\delta$, how many independent samples
$k$ are needed? The answer for every algorithm in this chapter is the
same family of bounds, derived from a single theorem about sums of
independent Bernoulli random variables. This section states that
theorem, sketches its proof via the moment-generating function, and
applies it to Miller-Rabin and to BPP amplification. The closing
remarks compare the bound to the variance-based confidence interval
of \Cref{sec:binomial-estimator}: same estimator, two different
tails.

%---------------------------------------------------------------------------
\subsection*{The multiplicative Chernoff bound}

Fix iid Bernoulli random variables $X_1, \ldots, X_n$ with success
probability $p$, and let $X = \sum_{i=1}^n X_i$ be their sum. Write
$\mu = \mathbb{E}[X] = np$ for the mean. The textbook variance bound
of \Cref{prop:binomial-variance} controls $\Pr[|X - \mu| \ge t]$
polynomially in $t$. The Chernoff bound improves the polynomial decay
to exponential decay in $\mu \epsilon^2$, at the price of an upper
tail expressed in multiplicative form $(1 + \epsilon) \mu$ rather than
additive form $\mu + t$.

\begin{theorem}[Multiplicative Chernoff bound]
\label{thm:chernoff}
Let $X_1, \ldots, X_n$ be independent Bernoulli random variables with
$\Pr[X_i = 1] = p$, and set $X = \sum_{i=1}^n X_i$ and $\mu = np$.
For every $\epsilon \in (0, 1]$,
\begin{equation}
\label{eq:chernoff-upper}
  \Pr[X \ge (1 + \epsilon) \mu] \;\le\; \exp\!\left(-\frac{\mu \epsilon^2}{3}\right),
\end{equation}
and symmetrically
\begin{equation}
\label{eq:chernoff-lower}
  \Pr[X \le (1 - \epsilon) \mu] \;\le\; \exp\!\left(-\frac{\mu \epsilon^2}{2}\right).
\end{equation}
\end{theorem}

\begin{proof}[Proof sketch]
The argument is the moment-generating-function (MGF) plus Markov
recipe, optimized over the free parameter. Fix $t > 0$. By Markov's
inequality applied to the random variable $e^{tX}$,
\[
  \Pr[X \ge a]
  \;=\; \Pr[e^{tX} \ge e^{ta}]
  \;\le\; e^{-ta} \, \mathbb{E}[e^{tX}].
\]
Independence of the $X_i$ factors the MGF of the sum into the product
of the per-variable MGFs:
\[
  \mathbb{E}[e^{tX}]
  \;=\; \prod_{i=1}^n \mathbb{E}[e^{tX_i}]
  \;=\; \prod_{i=1}^n \bigl(1 - p + p e^t\bigr)
  \;=\; \bigl(1 + p(e^t - 1)\bigr)^n.
\]
The pointwise inequality $1 + x \le e^x$ valid for all real $x$
(applied with $x = p(e^t - 1) \ge 0$) gives
\[
  \mathbb{E}[e^{tX}]
  \;\le\; \exp\!\bigl(np(e^t - 1)\bigr)
  \;=\; \exp\!\bigl(\mu(e^t - 1)\bigr).
\]
Substituting $a = (1 + \epsilon) \mu$ into the Markov step yields
\[
  \Pr[X \ge (1 + \epsilon) \mu]
  \;\le\; \exp\!\bigl(-t(1 + \epsilon) \mu + \mu(e^t - 1)\bigr).
\]
The right-hand exponent is minimized over $t > 0$ at $t = \ln(1 + \epsilon)$,
where it equals $-\mu \bigl((1 + \epsilon) \ln(1 + \epsilon) - \epsilon\bigr)$.
The function $\phi(\epsilon) := (1 + \epsilon) \ln(1 + \epsilon) - \epsilon$
satisfies $\phi(\epsilon) \ge \epsilon^2 / 3$ for $\epsilon \in (0, 1]$
(a Taylor-remainder estimate; the constant $1/3$ is the value that
makes the inequality tight at $\epsilon = 1$). Substituting gives
\eqref{eq:chernoff-upper}. The lower-tail bound \eqref{eq:chernoff-lower}
follows from the same recipe with $t < 0$ and the matching estimate
$\phi(-\epsilon) \ge \epsilon^2 / 2$.
\end{proof}

The two bounds are not symmetric in the constant: the upper tail has
exponent $\mu \epsilon^2 / 3$ and the lower tail has the tighter
$\mu \epsilon^2 / 2$. This asymmetry is a feature of the Bernoulli
case (large deviations above the mean are harder than below because
of the $\{0, 1\}$ ceiling) and matters in sample-complexity
calculations only when the dominant failure mode is on a specific
side. When both tails are bad and a two-sided bound is wanted, taking
the union of the two events and using the worse constant gives
\[
  \Pr[|X - \mu| \ge \epsilon \mu]
  \;\le\; 2 \exp\!\left(-\frac{\mu \epsilon^2}{3}\right),
\]
which is the form used in most of what follows.

%---------------------------------------------------------------------------
\subsection*{Application: Miller-Rabin round count}

The simplest application is also the cleanest, because Miller-Rabin
sits on the FNR-zero axis and the Chernoff bound is not needed at all
for the count. The per-round false-positive rate of
\Cref{thm:miller-rabin-formal} is $\fprate \le 1/4$, and the
$k$-round composition theorem reduces $\fprate^{(k)}$ to the product
of per-round rates: $\fprate^{(k)} \le (1/4)^k$. Choosing $k$ to
drive $\fprate^{(k)}$ below a target $\delta$ requires solving
$(1/4)^k \le \delta$, which gives
\[
  k \;\ge\; \log_4(1/\delta) \;=\; \frac{\log_2(1/\delta)}{2}.
\]
At $\delta = 10^{-6}$, $\log_2(10^6) \approx 19.93$, so $k \ge 10$
rounds suffice. At $\delta = 10^{-12}$, $k \ge 20$ rounds. At
$\delta = 2^{-128}$ (cryptographic-strength confidence), $k \ge 64$.
The bound is sharp up to the per-round constant $1/4$: doubling the
desired confidence in bits costs a constant number of additional
rounds, not a constant factor.

This calculation does not use \Cref{thm:chernoff} at all. The reason
is that Miller-Rabin's error is one-sided ($\fprate$ only, with
$\fnrate = 0$), so a single failure event collapses the chain. The
relevant bound is the chain rule of \Cref{thm:composition}, not a
concentration bound on a sum. Chernoff enters when the failure event
is a tail count rather than a single witness, which is the BPP case
below.

%---------------------------------------------------------------------------
\subsection*{Application: BPP amplification}

The BPP majority-vote amplification of \Cref{thm:bpp-amplification}
deferred its Chernoff calculation to this section. We can now supply
it. Let $A$ be a BPP algorithm with per-run accuracy bounded by
$1/2 + \epsilon_0$ on every input, where $\epsilon_0 > 0$ is the
margin above the $1/2$ threshold. (The conventional $2/3$ versus
$1/3$ of \Cref{def:bpp} corresponds to $\epsilon_0 = 1/6$.) Run $A$
on input $x$ for $k$ independent rounds, and report the majority of
the $k$ outputs. We bound the probability that the majority is wrong.

Let $X_i = 1$ if the $i$th run gives the correct answer, $0$
otherwise. By the BPP guarantee, $\Pr[X_i = 1] \ge 1/2 + \epsilon_0$
on every $x$, and by \Cref{ax:element-indep} the $X_i$ are mutually
independent. Let $X = \sum_{i=1}^k X_i$ be the number of correct
runs. The majority vote is wrong exactly when $X < k/2$. The mean of
$X$ is $\mu = k(1/2 + \epsilon_0)$, and the threshold $k/2$ sits
below $\mu$ by the deviation
\[
  \mu - k/2 \;=\; k \epsilon_0,
\]
which in multiplicative form is
\[
  k/2 \;=\; (1 - \epsilon')\mu
  \quad\text{with}\quad
  \epsilon' \;=\; \frac{\epsilon_0}{1/2 + \epsilon_0}.
\]
Applying the lower-tail bound \eqref{eq:chernoff-lower} gives
\[
  \Pr[X < k/2]
  \;\le\; \Pr[X \le (1 - \epsilon')\mu]
  \;\le\; \exp\!\left(-\frac{\mu (\epsilon')^2}{2}\right).
\]
Substituting $\mu = k(1/2 + \epsilon_0)$ and
$(\epsilon')^2 = \epsilon_0^2 / (1/2 + \epsilon_0)^2$, the right side
simplifies to
\[
  \exp\!\left(-\frac{k \epsilon_0^2}{2(1/2 + \epsilon_0)}\right)
  \;\le\; \exp\!\left(-\frac{2 k \epsilon_0^2}{3}\right)
\]
when $\epsilon_0 \le 1/2$. Loosening the constant to a clean
$\exp(-k \epsilon_0^2 / 3)$ and asking that this be at most $\delta$,
\[
  k \;\ge\; \frac{3 \ln(1/\delta)}{\epsilon_0^2}.
\]
For the conventional BPP gap $\epsilon_0 = 1/6$, this is
$k \ge 108 \ln(1/\delta)$, so $\delta = 10^{-6}$ requires
$k \approx 1490$ rounds. The constant $108$ looks large; in practice
the proof can be tightened by using \eqref{eq:chernoff-lower} with
its better $\mu \epsilon^2 / 2$ exponent throughout and the per-step
inequalities at their tightest, which trims the constant to single
digits. The qualitative point is what matters: BPP amplification is
$O((1/\epsilon_0^2) \log(1/\delta))$ rounds, exponential decay in $k$
rather than polynomial, and the exponent is the square of the margin
above $1/2$.

The cost asymmetry between BPP amplification and RP amplification is
worth naming. RP amplification (\Cref{thm:rp-amplification}) brings
$\fnrate$ from $1/2$ to $2^{-k}$ in $k$ rounds: a bit of confidence
per round. BPP amplification brings the same confidence in
$(3/\epsilon_0^2) k$ rounds, with a quadratic dependence on the
margin. A BPP algorithm with a $2/3$-versus-$1/3$ gap pays a
$\Theta(1/\epsilon_0^2) = \Theta(36)$ overhead per bit of confidence,
relative to a one-sided RP algorithm. This is the price of the
two-sided contract: when both errors can occur on the same coin
sequence, a single repetition is not enough; a vote across many
repetitions is needed, and votes concentrate exponentially in the
margin squared rather than exponentially in the count.

%---------------------------------------------------------------------------
\subsection*{Connection to the variance-based interval}

The Chernoff bound and the variance-based confidence interval of
\Cref{sec:binomial-estimator} are two tail bounds for the same
estimator. The binomial estimator $\hat{p} = X/n$ studied there had
variance $p(1-p)/n$ (\Cref{prop:binomial-variance}); Chebyshev's
inequality applied to that variance gives
\[
  \Pr\!\left[|\hat{p} - p| \ge \epsilon\right]
  \;\le\; \frac{p(1-p)}{n \epsilon^2}
  \;\le\; \frac{1}{4 n \epsilon^2},
\]
where the second step uses $p(1-p) \le 1/4$. The bound is polynomial
in $n \epsilon^2$: doubling $n$ halves the failure probability. The
Chernoff bound on the same estimator (rewriting $\hat{p} = X/n$ with
$X = \sum X_i$ and $\mu = np$, taking absolute deviation $\epsilon$
in the additive sense rather than multiplicative) gives
\[
  \Pr\!\left[|\hat{p} - p| \ge \epsilon\right]
  \;\le\; 2 \exp\!\left(-2 n \epsilon^2\right)
\]
by Hoeffding's inequality, the additive cousin of \Cref{thm:chernoff}
specialized to bounded $\{0,1\}$ random variables. The bound is
exponential in $n \epsilon^2$: doubling $n$ squares the failure
probability.

The two bounds describe the same estimator, but they answer different
questions at different cost. Chebyshev gives a bound that requires
only the variance, so it works for any distribution with the right
second moment; Chernoff requires the full MGF, so it is restricted to
distributions whose MGF is finite. The variance-based bound is what
chapter 7 §7.2 used to construct an unbiased point estimate with a
crisp Cramer-Rao certificate; the Chernoff bound is what chapter 11
uses to specify sample complexity. The two are not in competition;
they are the polynomial and the exponential tails of the same
binomial. When the question is ``what is the variance of $\hat{p}$,''
the chapter 7 answer suffices. When the question is ``how many
samples to drive the failure probability below $\delta$,'' the
chapter 11 answer is needed, because the polynomial bound forces
$n = \Theta(1/\delta)$ samples whereas the exponential bound asks
only $n = \Theta(\log(1/\delta))$. The gap between
$1/\delta$ and $\log(1/\delta)$ is the gap between a system that
needs $10^6$ samples for $\delta = 10^{-6}$ confidence and one that
needs $20$. This is why the Chernoff bound (and not Chebyshev) is the
workhorse of randomized-algorithm analysis.
